\section{Discussion} \label{sec:discussion}
% ============================================================================
\subsection{Comparison with Previous Works on Single-Pulse GRBs}
\label{sec:disc-compare}
Single-pulse bursts have long been the preferred laboratory for prompt-emission
spectroscopy, and our sample of $106$ bursts is, to our knowledge, the largest
uniformly-analyzed set to date. Table~\ref{tab:compare} summarizes the
comparison; we discuss it thematically.

\emph{The single-pulse lineage and spectral evolution.}
\citet{Ryde2004} and \citet{RydePeer2009} established that a single FRED-like
pulse exhibits the simplest spectral evolution and is well described by a
blackbody-plus-power-law, making single pulses the cleanest test of the prompt
mechanism. \citet{Lu2012} classified the \Ep\ evolution of bright single-pulse
Fermi bursts into hard-to-soft (HTS) and intensity-tracking (IT) patterns, a
scheme extended by \citet{Basak2013,Basak2014}. Our rise-phase classification
(HTS $63\%$, IT $25\%$; Section~\ref{sec:evol}) recovers the same two modes and
the same HTS-dominated balance as those small samples. The agreement on the
low-energy index is also notable: our median Band $\alpha=-0.71$ matches the IT
mean ($-0.68$) of \citet{Basak2014} and is close to the slow-cooled value
$-0.81\pm0.1$ of \citet{Burgess2014b}. On the synchrotron line of death the
comparison must be made carefully: \citet{Basak2014} report $\alpha>-2/3$ in
$63.4\%$ of their HTS \emph{bins} ($74.6\%$ excluding GRB~110721A), whereas our
like-for-like per-bin fraction is $44\%$ over a sample dominated by IT-era
decay bins, and $74\%$ ($74/100$) of our \emph{bursts} exceed the line in at
least one bin. The numbers are therefore comparable rather than identical, but
all the single-pulse studies---including \citet{Basak2013} and
\citet{Guiriec2010}---find the line of death frequently violated, which we now
establish at population scale.

\emph{The thermal / multi-component program.}
A sub-dominant photosphere has been advocated through blackbody composites
\citep{Ryde2004,Guiriec2010,Guiriec2011,Guiriec2015,Basak2014}. \citet{Basak2014}
find that a blackbody-plus-power-law peaks at $\sim3kT$, below the Band $\nu
F_\nu$ peak, and favor multi-blackbody models over Band; \citet{Guiriec2015}
require a cutoff-power-law$+$blackbody$+$power-law for two bright GRBs, with the
blackbody carrying only $1$--$10\%$ of the energy ($kT\sim40$~keV) and argue that
the hard $\alpha$ values are an artifact of fitting Band to a multi-component
spectrum. Our results support this picture at scale: the thermal flux fraction
is sub-dominant (median $0.006$ with interquartile range $0.001$--$0.10$, the
upper quartile matching their $1$--$10\%$ energy range), our decisive blackbody
temperatures ($\kT\approx3$--$70$~keV, median $28$) bracket their
$\sim$$40$~keV, and $91\%$ ($151/166$) of curvature-requiring bins are equally
or better described by adding a sub-dominant blackbody than by a second
synchrotron break.

\emph{The synchrotron / physical-model program.}
\citet{Burgess2014a} fit a physical synchrotron$+$blackbody model to six
single-pulse bursts and reported a strong thermal--non-thermal correlation
($\Ep\propto\kT^{\alpha}$, pooled $\rho\approx0.8$, $\alpha\sim1$--$2$); we test
this at $\sim$$18\times$ the sample size and recover it in $89\%$ ($17/19$) of
the bursts with enough decisive-blackbody bins to test, most cleanly in the
highest-S/N bursts (GRB~130427A: $\rho=0.93$, slope $0.88\pm0.07$), consistent
with empirical continua absorbing the thermal curvature in fainter bursts. The
double-break synchrotron picture \citep{Oganesyan2017,Ravasio2018,Ravasio2019}
predicts the 2SBPL we fit; in the brightest GRB ever, GRB~221009A,
\citet{Ravasio2023} find a genuine second low-energy break only in the
highest-S/N intervals---directly mirroring our result that an explicit second
break wins decisively in just $9\%$ ($15/166$) of curvature bins (a lower
limit), epitomized by GRB~110721A. This is expected: \citet{Ravasio2018} showed
that the $+$BB and two-break descriptions separate only in bright spectra with
a large peak-to-break ratio ($E_{\rm peak}/E_{\rm break}\gtrsim10$), precisely
the high-S/N regime. Finally, \citet{Mei2025} find that the
synchrotron model is acceptable for most bursts, that they lie in a
fast/intermediate-cooling regime ($\num\gtrsim\nuc$), and that the
spectral--energetics link is $\nuc$--$\Liso$ rather than the Yonetoku relation. We
independently support the absence of an $\Ep$--luminosity relation: the apparent
$\Ep$--intensity correlation (energy-flux $\rho=0.43$) collapses in photon flux
($\rho=0.15$), while the $\kT$--intensity correlation survives ($\rho=0.24$); and
our 2SBPL bins, with $\nuc<\num$ by construction, sit in the same
marginally-fast-cooling regime.

\emph{Large catalogs and the degeneracy lesson.}
Compared with the broad GBM catalogs \citep{Kaneko2006,Gruber2014,Yu2016} and the
short-GRB Bayesian catalog of \citet{Burgess2019}---in which $513/525$ short-burst
spectra prefer a cutoff power law with no thermal component---our six-model
analysis adds an explicit thermal-versus-two-break decomposition. Our central
caution echoes \citet{Basak2014}: no single empirical model is decisively
preferred over the next-best alternative in $92\%$ ($875/950$) of comparable
bins, so distinguishing the emission mechanism requires physical models
\citep{Burgess2014b} or higher signal-to-noise, not empirical goodness-of-fit
alone. Throughout, we verify our main results on the cleanest (Gold) subsample.

\subsection{The Two-Break Relation} \label{sec:disc-breaks}
The pooled $\num$--$\nuc$ correlation ($\rho=0.57$, $N=349$) can be decomposed
into two distinct statements: a \emph{between}-burst correlation (bursts whose
\nuc\ is higher also have higher \num\ on average) and a \emph{within}-burst
correlation (the two breaks move together in time inside an individual burst).
Both hold. Between bursts, $\rho=0.51$; within bursts, beyond the pooled
mean-subtracted residuals ($\rho=0.43$), the cluster-honest per-burst test
gives a median per-burst $\rho=0.50$, positive in $41/54$ bursts with
$\ge$$3$ double-break bins (sign-test $p=2\times10^{-4}$). The relation is
therefore not merely an inter-burst trend: harder bursts have both breaks at
higher energies, and within a burst the two breaks move together. Its slope,
however, depends on the estimator because the intrinsic scatter is large: the
\citet{DAgostini2005} fit gives $0.46\pm0.05$ (\ssc$=0.38$~dex), restricting to
bins with both breaks constrained gives $0.64\pm0.05$ ($\rho=0.67$, $N=253$),
and symmetric unweighted regression gives $1.05\pm0.06$; we therefore regard
the \emph{existence} of the correlation, not its slope, as the robust result.
It is also robust to the fit boundary: excluding the $26\%$ of \nuc\ values
within a factor of two of the 8~keV NaI edge leaves $\rho=0.55$ ($N=260$).
Because \nuc\ and \num\ are breaks of the same spectrum fit simultaneously,
part of the within-burst correlation can be induced by their joint estimation;
we note, however, that time-resolved \emph{physical} synchrotron fits of
GRB~180720B show the cooling and injection energies decreasing together
through the burst with $E_{\rm m}/E_{\rm c}\approx13$--$22$
\citep{Ronchi2020}---the same co-evolution, recovered with a model in which
the two breaks are not independently adjustable shape parameters. The
physically meaningful quantity is the ratio $\num/\nuc$ (the cooling
regime). Since the 2SBPL fit enforces $\nuc<\num$, the bins with a measurable
double break sit in the fast-cooling regime by construction; quantitatively,
the ratio distribution has a median $\num/\nuc\simeq6$ (interquartile
$3$--$13$), with $21\%$ ($75/349$) in the intermediate regime
$1<\num/\nuc<3$ of \citet{Mei2025}---close to the $\sim$$25\%$ they find in
their time-resolved sample, and well below the $\sim$$60\%$ of their
single-peak-bin sample. \citet{Mei2025} use this ratio as a cooling-regime
classifier rather than fitting a break--break relation; within the single burst GRB~160625B,
\citet{Ravasio2018} found $E_{\rm peak}$--$E_{\rm break}$ positively correlated
($\rho=0.61$), consistent in sign with our within-burst result. Our positive
$\num$--$\nuc$ correlation is, to our knowledge, the first such relation
reported for a large single-pulse \emph{sample}, and warrants confirmation
with physical synchrotron fits.

\subsection{Implication for the Radiation Process} \label{sec:disc-radiation}
The low-energy index peaking near $\alpha\approx-0.7$, between the
synchrotron limits but with a substantial fraction of bursts above them, is
consistent with marginally-fast/slow-cooled synchrotron with ongoing
acceleration in some bursts and a photospheric contribution in others
\citep[cf.][]{Basak2014,Burgess2014b}. The $\Ep$--$\kT$ slope can be mapped to
the outflow dynamics because the bulk Lorentz factor of an accelerating jet
grows with radius as $\Gamma\propto r^{\mu}$, with $\mu=1$ for a
baryon-dominated fireball and $\mu=1/3$ for a magnetically-accelerated jet;
the acceleration law sets how the photospheric temperature declines relative
to the non-thermal peak, so $\alpha\!\sim\!1$ indicates a baryon-dominated and
$\alpha\!\sim\!2$ a magnetically-dominated jet \citep{Burgess2014a}. The slope
we recover for the cleanest case, GRB~130427A
($\Ep\propto\kT^{0.88\pm0.07}$, \citealt{DAgostini2005} fit), is consistent
with the baryon-dominated value, but we stress that with empirical (not
physical) models this mapping is only indicative. The clearest two-break burst
in our sample, GRB~110721A (2SBPL-preferred, $F_{\rm BB}/F_{\rm
tot}\approx0.001$), is the counter-example in which the curvature is genuinely
a second synchrotron break rather than a thermal component.

% ============================================================================
